\section{Stage 9 --- Add optional SMB storage}

Integrate the NAS with AD before touching a workstation. Give each person a
private share, or a folder whose ACL grants that AD person and the storage
administrators access. Do not store a second copy of the person's password.
The workstation presents the current Kerberos identity to SMB.

\begin{gate}{The local profile is authoritative}
Windows must reach its desktop and Arch must reach its local home without the
NAS. Network storage is a convenience mounted after login. Never redirect
\texttt{Desktop}, \texttt{Documents}, \texttt{\%USERPROFILE\%}, or
\texttt{/home} during phase 1.
\end{gate}

\subsection*{Prove the server before configuring a client}

\begin{enumerate}
\item From AD DNS, resolve \placeholder{storage-host}. Record the returned
address and resolver. \textit{This catches split-DNS and stale-record errors.}
\item With an ordinary test identity, list the intended share. Confirm that a
different ordinary identity is denied. \textit{This proves the server ACL
instead of relying on a drive letter or desktop icon.}
\item Create a uniquely named test file, read it back, remove it, and record
the server-side owner. \textit{This detects guest fallback and identity-mapping
errors.}
\item Compare client, domain-controller, and NAS time. Keep observed skew below
the site's Kerberos limit. \textit{SMB may look like an ACL failure when the
ticket is actually too old or the clocks disagree.}
\end{enumerate}

\failurebranch{If the owner is guest, an unexpected numeric identity, or
another household user, stop. Correct the NAS domain join and ACL model before
automating any workstation mount. Do not compensate with a shared credential
or a permissive ACL.}

\subsection*{Windows: map after desktop login}

\begin{enumerate}
\item Log in online once as the ordinary domain test user. Run
\texttt{klist}; obtain a ticket before continuing.
\item In File Explorer, map the approved UNC path to the documented drive
letter. Select reconnect, but do not select alternate credentials and do not
save a password in Credential Manager.
\item Put the mapping action in a per-user logon task with an eight-second
timeout and no ``wait for network'' prerequisite. The task may report failure;
it may not hold the desktop open. Record the task name and exported,
secret-free definition.
\item Sign out and back in with the NAS reachable. Create a file, disconnect
the mapping, reconnect it, and read the same bytes. Record the path, SHA-256
digest, elapsed login time, and effective AD owner.
\end{enumerate}

\proofpair{Get-SmbMapping}{The approved drive points to the approved UNC path,
uses the signed-in domain identity, and reports no guest mapping.}
\proofpair{cmd /c ``net use''}{No obsolete drive letter, alternate NAS name,
or unexpected cached connection is present.}

\subsection*{Arch: mount in the user session}

Install the repository-declared SMB and desktop virtual-filesystem packages.
Use the desktop's user-session SMB mount, such as GVfs, so no network
filesystem participates in boot or local-home creation.

\begin{enumerate}
\item Log in online as the ordinary domain test user. Run \texttt{klist}; fix
SSSD or Kerberos first if no current ticket exists.
\item Open \texttt{smb://\placeholder{storage-host}/\placeholder{user-share}}
from the desktop file manager. Choose the current domain identity and decline
permanent password storage.
\item If automatic connection is wanted, create a user-session service that
runs the mount after the graphical session starts, has an eight-second
timeout, and is not required by the session target. Never add the SMB path to
\texttt{/etc/fstab}, the boot target, or the local-home unit.
\item Create, hash, unmount, remount, and read the same test file. Record the
Kerberos principal and effective server-side owner.
\end{enumerate}

\proofpair{gio mount -l}{The approved SMB location appears only when reachable;
the local home remains an ordinary local filesystem.}
\proofpair{systemctl --user --failed}{A failed optional mount may be visible,
but no login, home, or graphical-session unit depends on it.}

\subsection*{Run the three-state failure drill on both systems}

\begin{tabularx}{\linewidth}{@{}p{.17\linewidth}p{.31\linewidth}L@{}}
\toprule
\textbf{State} & \textbf{Action and measurement} & \textbf{Pass boundary}\\
\midrule
Available & Login, create and hash a file; unmount and remount & Same bytes,
same AD owner, no password prompt\\
Unreachable & Block or disconnect only the NAS; cold-login with a stopwatch &
Local desktop/home opens within the declared limit; bounded warning only\\
Denied & Keep the host reachable but remove the test user's share access &
Explicit access-denied result; local login works; no guest or alternate user\\
\bottomrule
\end{tabularx}

\askbefore{Did each test begin from a cold sign-in rather than an already
connected session? Was the NAS itself unreachable in the second test, but
reachable in the third? Did either client silently reuse an old connection or
credential?}

\failurebranch{If login waits on storage, remove the mapping from every boot,
profile, shell-startup and required-session dependency, then repeat all three
states. If access denied becomes access granted after a prompt, remove the
saved credential and inspect the server ACL.}

\section{Stage 10 --- Add restricted managed Wi-Fi}

Create a dedicated managed-workstation SSID and VLAN. Prefer one PPSK per
workstation so one lost laptop can be removed without changing every machine.
Retrieve the selected key from the private secret store only while minting the
machine. Phase-1 disks are unencrypted; assume physical access can recover any
stored wireless credential.

\begin{center}
\begin{tikzpicture}[x=1cm,y=1cm,font=\scriptsize]
  \draw[wf pencil,rounded corners=1pt] (-5.6,-.10) rectangle (-3.35,1.22);
  \draw[wf faint] (-5.42,.08) rectangle (-3.53,1.04);
  \node[align=center] at (-4.48,.56) {managed\\workstation};
  \draw[wf flow] (-3.35,.56)--(-1.63,.56) node[midway,above]{PPSK};
  \draw[wf pencil] (-1.63,.18)--(-1.15,.56)--(-1.63,.94);
  \draw[wf pencil] (-1.35,.05)--(-.71,.56)--(-1.35,1.07);
  \draw[wf pencil] (-1.03,-.09)--(-.20,.56)--(-1.03,1.21);
  \node[below] at (-.92,-.13) {restricted SSID};
  \draw[wf flow] (-.20,.56)--(1.10,.56);
  \draw[wf pencil] (1.10,-.10) rectangle (2.65,1.22);
  \node[align=center] at (1.88,.56) {AD/DNS\\updates\\approved SMB};
  \draw[wf deny] (3.12,-.10)--(3.12,1.22);
  \draw[wf pencil] (2.91,1.06)--(3.33,.06);
  \draw[wf pencil] (2.91,.06)--(3.33,1.06);
  \draw[wf faint] (3.72,-.10) rectangle (5.58,1.22);
  \node[align=center] at (4.65,.56) {management\\household peers};
\end{tikzpicture}
\wfcaption{A stored credential reaches only the services a managed workstation needs.}
\end{center}

\subsection*{Prove UniFi policy before storing a key}

Using a temporary test client, verify association, DHCP lease, gateway, AD
DNS, time, domain ports, approved SMB, and both operating systems' update
destinations. Then verify denials to gateway administration, Controller
administration, infrastructure management, household-client, peer-client, and
IoT destinations. Save a small source/destination/port/result table.

\failurebranch{If an allow fails, inspect association, DHCP, DNS, time, then
the single destination and port. If a deny fails, quarantine the SSID. Do not
store its credential on a workstation until both directions match policy.}

\subsection*{Windows: install one machine profile}

\begin{enumerate}
\item Create a WLAN profile for only the managed SSID and its assigned PPSK.
Import it for all users from the local console. Do not add household,
management, or provisioning SSIDs.
\item Delete the plaintext import file immediately after the profile appears.
Check Git status, the build log, answer files and shell history for the SSID
key; store only the key's secret-manager reference in evidence.
\item Disconnect Ethernet, restart, and log in with a previously cached
ordinary account. Confirm automatic association and the expected managed-VLAN
lease.
\item Forget or disable the profile, restart, and prove that Windows does not
fall back to another known private network. Restore only the approved profile.
\end{enumerate}

\proofpair{netsh wlan show profiles}{Exactly the approved preload set appears;
captured output contains no key material.}
\proofpair{Get-NetConnectionProfile}{The active wireless interface has the
expected network category and no second active route into a private zone.}

\subsection*{Arch: install one root-owned NetworkManager profile}

\begin{enumerate}
\item Create one system connection for the managed SSID with NetworkManager,
binding it to the workstation PPSK and enabling autoconnect. Avoid typing the
key in a command that will be retained in shell history.
\item Confirm its file under
\texttt{/etc/NetworkManager/system-connections/} is owned by root, mode 600,
and excluded from logs, artifacts, backups without secret encryption, and Git.
\item Disconnect Ethernet, restart NetworkManager or reboot, and log in with a
previously cached ordinary account. Confirm association, expected lease, AD
DNS, time, and the default route.
\item Set the connection down and prove that no household or management SSID
autoconnects. Bring back only the approved profile.
\end{enumerate}

\proofpair{nmcli -f NAME,TYPE,AUTOCONNECT connection show}{Only approved
wireless profiles autoconnect; do not include secret fields in evidence.}
\proofpair{stat -c '\%U \%G \%a \%n'
/etc/NetworkManager/system-connections/*}{Every stored system profile is
root-owned and mode 600.}

\subsection*{Wireless-only acceptance}

Repeat AD DNS SRV lookup, clock check, ordinary online login, administrator
authorization, OS update check and optional-SMB three-state drill with Ethernet
physically disconnected. Measure association-to-lease and login time. Scan
each forbidden zone from only its documented small CIDR; do not substitute a
wide private-range scan.

\askbefore{Does revoking this workstation's PPSK leave every other workstation
connected? Does the revoked machine fail to associate after its cached radio
state is cleared? Can it reach all required update endpoints but no management
page or peer workstation?}

\failurebranch{If wireless onboarding fails, continue over the deployment
cable. Never inject the household Wi-Fi key or broaden the VLAN to make the
test pass. If a required vendor update cannot be bounded yet, record that
exception and keep the workstation a pilot.}
